% English arXiv version of free-research-012 (run-012), translated from the run's Japanese paper.md.
% Tables live in tables/*.tex and are generated by ../make_tables_tex.py from the run's analysis outputs.
% \fillin{...} = facts only the human author can supply; leave them red until supplied.
% Keep the source ASCII-only (arXiv builds with pdfLaTeX).
\documentclass[11pt]{article}

\usepackage[T1]{fontenc}
\usepackage[utf8]{inputenc}
\usepackage{lmodern}
\usepackage[margin=1in]{geometry}
\usepackage{microtype}
\usepackage{amsmath}
\usepackage{booktabs}
\usepackage{graphicx}
\usepackage{xcolor}
\usepackage{enumitem}
\usepackage[font=small,labelfont=bf]{caption}
\usepackage[hidelinks]{hyperref}
\usepackage{url}
\usepackage{adjustbox}
\usepackage{float}
\usepackage{rotating}
\emergencystretch=2em
\setlength{\tabcolsep}{4pt}

\newcommand{\fillin}[1]{\textcolor{red}{[#1]}}        % to be completed by the human author before submission
\newcommand{\hdl}[1]{\begin{tabular}[b]{@{}l@{}}#1\end{tabular}}  % multi-line table header, left
\newcommand{\hdr}[1]{\begin{tabular}[b]{@{}r@{}}#1\end{tabular}}  % multi-line table header, right
\newcommand{\gentable}[1]{\adjustbox{max width=\linewidth}{\input{tables/#1.tex}}}

\title{Forecasters and the Freshness of Published Information in Wastewater Surveillance\\[2pt]\large A Comparison Based on Actual Publication Histories}

\author{Ko Fujimura\thanks{\texttt{fujimura@otsuma.ac.jp}. This study was conceived, executed and written autonomously by an AI research agent (OpenAI Codex with the model \texttt{gpt-6.1-sol}; see the AI Declaration) in one session on 18--19 October 2026; the human author's role is described in the AI Declaration at the end of the paper. Apart from self-checks by the same AI and one review round by another AI agent, the work has not been peer reviewed or independently replicated.}\\[2pt]
{\normalsize Otsuma Women's University, Faculty of Social Information Studies}}
\date{}

\begin{document}
\maketitle

\begin{abstract}
In wastewater surveillance, publication delays and revisions of past values affect judgements about infection trends. Using official Git timestamps as a proxy for publication time, this study compared pretrained general-purpose time-series forecasters fed older published inputs with classical forecasters fed the latest inputs. In the pre-planned German national analysis, the weighted interval score (WIS) ratio with input published one week earlier was 1.039 [0.915, 1.163] for Chronos-Bolt and 1.023 [0.932, 1.155] for TimesFM (latest-input three-method classical ensemble = 1; smaller is better). This average includes gaps in the publication history and shows neither equivalence nor non-inferiority. The one-week-earlier input was identical to the latest input in 75/338 and 22/262 comparisons; where the latest input had newer observation weeks, the ratios were 1.147 and 1.064. Relative to the same method with the latest input, one-week-earlier input gave ratios of 1.233 (classical ensemble), 1.133 (Chronos-Bolt) and 1.200 (TimesFM), so differences between forecasters and differences due to information updates must be judged together. Results also differed by pathogen. Doublings were limited to 6 cases per target position, so detection performance could not be concluded. In an exploratory post-review analysis of pairs whose latest-reference horizon was at most two weeks, the Chronos-Bolt ratio with one-week-earlier input was 1.127 [1.034, 1.222]. Practitioners working at short horizons should not replace the latest input with older general-purpose forecasts solely on the basis of averages that include publication pauses. Causal effects of faster measurement or publication, costs and effects on clinical alerts were not examined. The study was conducted end-to-end by an AI agent (see the AI Declaration).
\end{abstract}

\section{Introduction}

Wastewater surveillance is a means of continuously tracking the trends of pathogens shed by a population, not only by individuals who have been tested. What practice needs is not only to explain the curve that has been tidied up afterwards, but also to judge, from the information available now, whether concentrations are rising and how they will change in the next week. Sampling, measurement, aggregation and publication, however, take time. When infection dynamics are estimated from concentrations, biological time lags also differ from operational delays such as those of laboratories~\cite{lison2026}.

Surveillance staff and forecasting researchers working on this problem face a design decision: whether to change the forecaster or to use fresher measurement and publication information. Pretrained time-series models can produce forecasts without additional training of weights for each target. An evaluation that merely truncates a series finalised later at each time point, however, differs from an evaluation that uses the values that had actually been published at that time. Moreover, when publication is not weekly, the version published one week ago and the current version can have the same observation end. Before a model difference is read as ``compensating for so many weeks of information'', it is therefore necessary to check which inputs and reference values actually changed.

The research question is: \emph{how well can pretrained general-purpose time-series forecasts that take actually published information from one to two weeks earlier reproduce the wastewater concentrations of the same week and their rises, compared with classical forecasts that take the latest information, and how does the comparison change when the addition of observation weeks is distinguished from revisions of existing values?} The study compares the accuracy of indicators. It does not estimate the trade-off between the costs of testing and of running models, or the effects on medical decisions.

\paragraph{Summary of findings.}
\begin{enumerate}[leftmargin=*,itemsep=2pt]
\item In the pre-planned German national analysis, the WIS ratio with one-week-earlier input, relative to the fixed three-method classical ensemble with the latest input, was 1.039 [0.915, 1.163] for Chronos-Bolt and 1.023 [0.932, 1.155] for TimesFM (Table~\ref{tab:primary}). This average includes gaps in the publication history and shows neither equivalence nor non-inferiority.
\item The one-week-earlier input was identical to the latest input in 75 of 338 (Chronos-Bolt) and 22 of 262 (TimesFM) comparisons; in the group whose latest input had newer observation weeks, the ratios were 1.147 and 1.064 (Table~\ref{tab:information}).
\item Relative to the same method with the latest input, one-week-earlier input gave ratios of 1.233 for the classical ensemble, 1.133 for Chronos-Bolt and 1.200 for TimesFM (Table~\ref{tab:ownage}). Results also differed between pathogens (Table~\ref{tab:pathogens}).
\item Doublings were limited to six cases per target position, so detection performance could not be concluded (Table~\ref{tab:alerts}).
\item In an exploratory analysis added after review, among pairs whose latest-reference horizon was at most two weeks, the ratio of Chronos-Bolt with one-week-earlier input was 1.127 [1.034, 1.222] (Table~\ref{tab:horizons}).
\item The exploratory extensions (German sites, New Zealand, stronger classical references, pooled learning, longer histories) did not show a uniform relation (Section~\ref{sec:explore} and Appendix~\ref{app:results}).
\end{enumerate}

\section{Related Work and Differences}

Neither forecasting with wastewater information nor the handling of reporting delays is a new idea. Radermacher et al.~\cite{radermacher2025} evaluated short-term forecasts of hospitalisations in Germany using wastewater information and, in their setting, found no clear benefit from adding wastewater information to forecasts based on hospitalisation history alone. That paper assumes that the data up to the time of use are complete. EpiSewer estimates infection dynamics from wastewater and also evaluates future measured concentrations~\cite{lison2026}. EpiFlow forecasts disease indicators from wastewater and includes simulations that delay reporting~\cite{adiga2026}.

Nowcasting that exploits additions and revisions of reports has a long history, including Bayesian smoothing of reporting-delay triangles~\cite{mcgough2020}. There are also studies that forecast several surveillance indicators directly from revision histories, and their targets are not limited to integer counts~\cite{tang2025}. The official CDC implementation also performs real-time forecasting of wastewater and hospitalisations and retrospective evaluation with saved versions. Its official documentation records that versions of the wastewater inputs were saved every week, so this study does not regard the use of actual publication histories as novel in itself~\cite{cdcpipeline}. The corresponding 2026 preprint was found by search, but direct retrieval of its full text returned HTTP 403, so it is not used as evidence for detailed results. The perspective of separating the timeliness and the quality of information also exists in real-time forecasting of economic statistics~\cite{giannone2005}.

The difference of this study is not the invention of a new delay-correction model or score. Within the range searched, no evaluation was found that directly compares pretrained time-series forecasts with the latest classical forecasts on actual wastewater publication histories and that separates, together, \emph{identical inputs, updates of existing history, additions of observation weeks, and the scale of the later reference}. This is not an assertion that no precedent exists. The extensions to sites and another country likewise examine the scope of the same question.

Evaluation with the versions available at publication time has a long tradition. Croushore and Stark~\cite{croushore2001} built a real-time data set for economic statistics and compared AR forecasts against different reference values. The supplement ``Finalized Versus Vintage Data'' of McDonald et al.~\cite{mcdonald2021} shows that evaluating COVID indicators with finalised data yields performance that differs by method from evaluation with actually published inputs. This study does not present this evaluation principle as a new method.

Applying pretrained time-series models to infectious diseases is also established. EpiCastBench~\cite{panja2026} compared models including TimesFM and Chronos-2 on several infectious-disease datasets, reporting low average errors and that classical machine learning is also competitive at short horizons. Kalahasti et al.~\cite{kalahasti2025} is also related, but in this study only the abstract on the authors' publication page was checked, and retrieval of the full text returned HTTP 403, so it is not used as evidence for specific performance. The ratios of this study with the latest input are not positioned as the first evidence that general-purpose models can be used for infectious-disease forecasting.

For forecasting wastewater indicators themselves, Alhassan et al.~\cite{alhassan2025} compared several classical models on the US National Wastewater Surveillance System (NWSS) and report differences by horizon and region. The difference of this study is not forecasting future wastewater concentrations or the existence of horizon dependence. It is comparing the combination of the freshness of actually published inputs and general-purpose forecasts with the latest classical forecasts, and making explicit the update conditions that make up the average.

Hu et al.~\cite{hu2026}, identified in an additional search, also propose a forecaster that explicitly represents reporting frequency and revisions, and compare it with TimesFM and other models on case and hospitalisation data. The emphasis on the publication process and the mismatch between general-purpose models and surveillance data are thus also existing topics. This study does not discover these general points; it examines how the comparison between old and latest information changes with actually published wastewater inputs.

A re-search before submission also identified VINTAGE-TS~\cite{ahmad2026}. It separates publication time and observation time, defines first and matured references, and proposes revision-aware handling for general-purpose time-series models. Its executed numerical evaluation consists of linear methods on synthetic data and the like, and it states that a comparison of real economic data with Chronos-2 has not been executed. The difference of this study lies not in this general evaluation principle but only in comparing the losses of the latest classical methods and old general-purpose forecasts on actual wastewater publication histories.

\section{Data and Definitions of Time}\label{sec:data}

\subsection{Main analysis}

The changed versions of the national aggregate were frozen from the official public Git history of Germany's AMELAG wastewater surveillance. The frozen publication history has 120 versions, the last being commit \texttt{6586326ea292f4bfa243d8444d1efc8f4fef772d}. The data are published by the Robert Koch Institute and the German Environment Agency (Umweltbundesamt) under CC BY 4.0~\cite{amelagdata}. Sources, retrieval times, hashes and documents are recorded in the replication package.

The targets are measured concentrations of SARS-CoV-2, influenza A and RSV A/B. For SARS-CoV-2 the published normalised indicator is used, and for the other two series the non-normalised measured indicator. RSV A/B is not chained to the later RSV A+B as if they had the same definition. Columns of smoothed curves are not treated as forecasts. National week labels are Wednesday-ending weeks from Thursday to Wednesday; ineligible published versions whose dates fall on a different weekday fell back to the immediately preceding eligible version. This rule was decided in an input audit before forecasting, and the excluded versions were preserved.

The official description of the frozen latest version states that the weekly report uses the non-normalised \texttt{viruslast} and that normalisation did not improve data quality. It also describes the national aggregation, which uses series means of sites and laboratories. The existence of this processing and of revisions is known and is not a finding of this study. To examine the dependence on the normalised indicator chosen for the main analysis, an additional comparison with non-normalised SARS-CoV-2 was planned~\cite{amelagreadme}.

The updates of the saved national aggregate file had gaps of 56 days from 2025-06-18 to 2025-08-13, 43 days from 2025-09-03 to 2025-10-16, and 15 days from 2025-12-23 to 2026-01-07. This is the publication history of that file; it was not confirmed whether sampling or measurement, or other distribution channels, stopped.

Issues for calibration started on 2025-01-03, and the evaluation covered 58 weeks from 2025-07-04 to 2026-08-07. Relative to the Friday issue time, the latest versions published by that time, by one week earlier and by two weeks earlier are used as inputs. The targets are the preceding Wednesday and the next Wednesday. ``One week earlier'' refers to the cutoff for usable published information and does not mean that the last measurement date is necessarily one week older.

The main reference value is the measurement in the latest version available by the end of the day 56 days after the label date of the target week. It is neither the true viral load nor a final definitive value. References at 28 days, at 84 days and from the frozen final published version are used in sensitivity analyses. Values that have not passed the 56-day cutoff are not used for learning from past performance. Because Git committer timestamps were used as a proxy for publication time, there is no guarantee that they coincide exactly with independently confirmed release times.

\subsection{Exploratory generalisation}

In Germany, 71 sites and 157 series were fixed from public information before testing, on the basis of population size, observation history and recent observation frequency. Only positive measurements are used; several values on the same day are averaged on the log scale, and the daily values within a week are then averaged on the log scale. Rows that are only smoothed are not treated as measurements. In New Zealand, one national and five regional series from the official history~\cite{nzdocs}, published under CC BY 4.0, were used.\footnote{Attribution required by the data's terms of use: COVID-19 Data Repository by the New Zealand Institute for Public Health and Forensic Science Ltd~\cite{nzdata}. The repository asks academic publications to cite~\cite{hewitt2022}.} Six versions in the whole history contained malformed date rows or similar defects; they were not repaired, and the immediately preceding eligible version was used. Issues are on Mondays for Sunday-ending weeks. Unlike the German normalised indicator, the unit is genome copies per person per day. Both extensions are exploratory checks planned after the main results.

In the saved New Zealand publication history of the test period there were 30 publication days; of the 29 consecutive intervals, 16 were 14 days, so publication was roughly fortnightly. Thursday was the most frequent weekday (24 days). Even with a one-week-earlier cutoff, cases without an update and cases in which several weeks are added at the next publication are mixed, so it does not mean the same observation difference as the German one-week-earlier input.

No personal information or patient data were collected. Sites are published aggregation units, and no linkage with small numbers of households or individuals was attempted. The results do not show the superiority or inferiority of individual facilities. The hash labels of sites and regions are display identifiers, not a guarantee of anonymisation. Because definitions, sample density, normalisation and publication frequency differ, the results cannot be read as rankings of countries or organisations.

\section{Methods}\label{sec:methods}

\subsection{Forecasts and comparison conditions}

Forecasts are made on the common logarithm of concentration. Histories have at most 104 weeks and at least 16 measured values, and internal missing values in a series are forward-filled with past values. No interpolation from the future is made. Conditions whose horizon to the target exceeds 16 weeks are kept as out of range. Conditions that forecast a target week already contained in the input are counted separately.

The public weights of Chronos-Bolt base~\cite{chronosbolt} and TimesFM 2.5~\cite{timesfm25} were run on a local CPU. The weights were not trained further, and no target region or pathogen names were given as input. What is passed to the models is a column of log values; dates are also included when inputs are identified. Saved histories are in double precision, numeric input to the general-purpose models is float32, and the classical methods use double precision. The publication history of the weight files is recorded in the replication package: Chronos-Bolt was published before the start of the test period, and TimesFM is limited to evaluations from 2025-10-03 onwards. TimesFM forecasts before its publication are used for calibration calculations and checks of numerical consistency, and are not interpreted as the performance that could have been obtained in operation at that time. The full content of the pretraining data could not be checked, and no claim is made that it does not overlap with past wastewater data.

The nine quantiles of each forecast were re-sorted in ascending order, and the resulting median quantile was used as the point forecast. This rule was fixed before forecasting. Even when numeric columns are identical, a measured value in a later version can coincide with an earlier imputed value, so identical input is not called identical measurement provenance. For strata of imputed proportions, the provenance of each published version is used.

The pre-planned classical reference is a fixed ensemble that averages the quantiles of three methods with equal weights: non-seasonal AutoARIMA, AutoETS with additive errors, and a forecast that applies damping to the slope of the last eight points. In the tables it is labelled ``Fixed 3-method''. A last-value-carried-forward forecast is also computed. The ARIMA orders, differencing and selection criteria, the ETS model set and the variance estimation were fixed in the plan and code. The background of these standard models follows Hyndman and Athanasopoulos~\cite{fpp3}.

After the main results had been seen, a fixed six-method ensemble that adds AutoTheta, a forecast using the same week of the previous year and ARIMA with harmonic regression was added, together with a method that selects the best classical method from the most recent 26 cases whose 56-day reference cutoff had passed by the time of use. At least 12 cases are required, and the best method is not chosen from future test performance. These are exploratory comparisons; the tables label them ``Fixed 6-method'' and ``Past selection''. For all methods, a condition that corrects quantiles using available past residuals is also computed. No coverage guarantee from the correction is assumed.

For the additional comparisons for sites and New Zealand, gradient-boosted quantile regression trained per country on several pooled past series was also used. The features are past level, slope, variability, horizon, pathogen category and the annual cycle of the target date; site IDs and future values are not included. The model is trained only on cases whose 56-day reference cutoff passed before each issue time, and hyperparameters are not tuned on the test. This method is not a comparison that equalises the amount of information with the general-purpose models. Settings, features, training rows and models were saved~\cite{sklearnhgb}.

To examine the dependence on the interval-correction method, a moving-window correction based on the nonconformity score of conformalized quantile regression was also added. The endpoints are corrected on the basis of past exceedances of the central interval, and the median is kept fixed. Because exchangeability of time series cannot be assumed, the coverage guarantee of that method is not applied to this study. Details and all windows are given in Appendix~\ref{app:calibration}~\cite{romano2019}.

\subsection{Evaluation and uncertainty}

The main metric is the weighted interval score (WIS); auxiliary metrics are the mean absolute error (MAE) of the median and the coverage of the 80\% prediction interval. WIS is twice the mean pinball loss over the nine quantiles from 0.1 to 0.9, and it was cross-checked against the definition based on central intervals. All values are on the log-concentration scale, and smaller is better. The score definitions are existing ones~\cite{bracher2021,bracher2022}.

A ratio is the ratio of the mean losses of the candidate and the reference over the same targets that can be evaluated for both. A value below 1 means the candidate's loss is smaller. The main intervals are 95\% intervals from moving-block resampling with 8-week blocks; 4- and 12-week blocks are also saved. In each replicate, the series and target weeks belonging to the same issue week were resampled together. The number of replicates is 2,000, and the seed is 12012. The number of series is not counted as the number of independent epidemics. The nominal intervals in the main tables are not adjusted for multiple comparisons, and an interval that straddles 1 is not taken as proof of equivalence. No non-inferiority margin was set in advance either.

A rise is defined as a concentration at least twice that of two weeks earlier in the same reference version. Alerts are generated from changes of the point forecast, with the 90th percentile of past non-rise cases as the threshold. The intended false-alert rate is 10\%, but there is no guarantee that this rate holds in the test. No rise probability is constructed by assuming dependence between the two time points of the forecast distribution. The receiver operating characteristic (ROC) curve for doubling with the original methods is a supplement to the pre-registered plan; ROC curves and areas under the curve (AUC) for other rise ratios and for the additional methods are kept separate as exploratory analyses. The ROC describes discrimination looking back over the whole test period and differs from the performance of thresholds available in the past. A rise is a mechanical change of a published wastewater indicator and not the ground truth of a clinical alert.

\subsection{Checks of information updates and of scale}

For candidates and references, three kinds of input are distinguished: inputs whose saved dates and log values are identical; inputs with the same observation end but changed values or similar differences; and inputs whose measurements extend to newer weeks. ``Identical input'' is defined on the saved history and does not mean that all metadata of the publication are identical. Agreement after float32 conversion and histories with close values are examined separately in another exploratory check.

In an exploratory decomposition, an old actually published input A, an input B obtained by truncating the latest published version at the observation end of A, and the latest input C are compared. The difference from A to B includes revisions and acceptance of existing weeks and observations that were added late; the difference from B to C includes the extension of the observation end. This is a descriptive comparison in a fixed order, not the causal effect of an intervention in the publication process. An analysis that measures changes from a common past base week in both forecasts and references was also added, to distinguish revisions of absolute levels from the forecasting of trends. A condition that truncates the later final version at past observation ends is also computed; because it contains future revisions, it serves only for diagnosis and is not compared as an available forecaster.

After review, an exploratory analysis was added that splits pairs by the horizon of the latest three-method reference into at most two weeks and three weeks or more. The boundary was proposed by the reviewer after seeing the results. The same evaluation calendar as the original, with 8-week blocks and 2,000 replicates, was used, and weeks outside a group were kept with zero cases. The horizon groups are not treated as definitive labels of regular publication versus pauses, and the analysis is neither a causal effect of publication nor a test adjusted for multiple comparisons.

\section{Results}\label{sec:results}

\subsection{Pre-planned main results}

There were 348 intended test conditions with the latest actually published information. The main common comparison had 338 cases; 10 cases lacked the main reference, and 0 conditions contained a known target week in the input.

\begin{table}[tbp]\centering\small
\caption{WIS ratios relative to the fixed three-method ensemble with the latest input (German national analysis, pre-planned). The classical ensemble's own older inputs are listed as reference rows. Brackets: 95\% intervals from 8-week blocks. TimesFM covers only the period after its publication, so its sample size differs from Chronos-Bolt. Within a row, the candidate and the reference are compared on the same targets.}\label{tab:primary}
\gentable{tab_primary}
\end{table}

With the latest information as input, the ratio was 0.917 [0.812, 1.000] for Chronos-Bolt and 0.852 [0.789, 0.934] for TimesFM. The ratios with older published information are shown in Table~\ref{tab:primary}; they give no basis for converting the model improvement into a fixed ``number of weeks of information''. In the group in which the one-week-earlier information differed from the latest by newer observation weeks, the Chronos-Bolt ratio was 1.147 [1.058, 1.246]. This grouping, however, also depends on the publication process and the epidemic period, so it is not interpreted as a causal effect of adding observations.

\begin{table}[tbp]\centering\small
\caption{WIS ratios by input relation when the one-week-earlier published information is used. These groups depend on the data-update process and are not randomly assigned conditions. The intervals of groups with few issue weeks, especially the identical-input group of TimesFM, are for reference only and are not taken to rest on enough independent epidemic episodes.}\label{tab:information}
\gentable{tab_information}
\end{table}

\begin{table}[tbp]\centering\small
\caption{Pre-planned breakdown by pathogen. For each method, pathogen and period, inputs from 0, 1 and 2 weeks earlier are compared with the latest fixed three-method ensemble on the same targets. For TimesFM, the full test period also includes only the period after its publication. Intervals are not adjusted for multiple comparisons, and the point estimates leaning towards a disadvantage for SARS-CoV-2 are not replaced by the average over all pathogens.}\label{tab:pathogens}
\gentable{tab_pathogens}
\end{table}

\begin{table}[tbp]\centering\small
\caption{Exploratory comparison by horizon, added after review. Each column takes the latest fixed three-method ensemble as 1. The horizon is the number of weeks from the observation end of the latest reference to the target week. Issue weeks can overlap between groups, so they are not added up.}\label{tab:horizons}
\gentable{tab_horizons}
\end{table}

In the exploratory post-review analysis, among the 253 pairs whose latest-reference horizon was at most two weeks, the ratio with one-week-earlier input was 1.127 [1.034, 1.222] for Chronos-Bolt and 1.054 [0.954, 1.170] for TimesFM (Table~\ref{tab:horizons}). The groups are a division by horizon and do not completely separate regular publication from pauses. The group with three weeks or more also includes publication delays of a few days and next-week targets, and the same issue week can enter both groups. The overall average of each method must be distinguished from its performance when operating at short horizons.

\begin{table}[tbp]\centering\small
\caption{Post-review exploratory summary that compares only the freshness of the published input within the same method. Intervals are not adjusted for multiple comparisons. Even for the same statistic, the interval endpoints differ slightly from Table~\ref{tab:primary} because resampling used a different random stream. Horizons and calendar periods were not randomly assigned. TimesFM covers only the period after its publication, so the evaluation periods of the three rows are not the same. The order of the ratios across rows is not to be read as a ranking of methods.}\label{tab:ownage}
\gentable{tab_own_age}
\end{table}

In pairs that differ from the latest input by newer observation weeks, the WIS ratio of the one-week-earlier input, with the same method's latest input as reference, was 1.165 [1.112, 1.235] for Chronos-Bolt, 1.233 [1.170, 1.313] for TimesFM and 1.315 [1.234, 1.421] for the three-method ensemble. This is a comparison conditional on publication timing, not the effect of an intervention that speeds up measurement or publication.\footnote{The original manuscript also plotted the national ratios against information age for the three classical references. The plotted values are given in Table~\ref{tab:primary} and Table~\ref{tab:extended}.}

\subsection{Detection of rises}

\begin{table}[tbp]\centering\small
\caption{Detection of doublings (a rise to at least twice the value of two weeks earlier) by alerts from point forecasts, with one-week-earlier input, compared with the latest fixed three-method ensemble on the same cases. The offset is the target date minus the issue date.}\label{tab:alerts}
\gentable{tab_alerts}
\end{table}

The doublings were limited to 6 and 6 cases for the two target positions (Table~\ref{tab:alerts}). Because the same wave of rises overlaps across target positions and adjacent weeks, these are not added up as independent clinical alert events. A smaller concentration error does not necessarily improve the false-alert or miss rate. Some conditions exceeded the intended false-alert rate, and clinical usefulness cannot be concluded from this sample.

\begin{table}[tbp]\centering\small
\caption{Point-forecast error and intervals for the main one-week-earlier input and the latest reference, aligned on the same targets. Interval widths are on the common-log concentration scale. Because each model is evaluated only after its publication, this is not a direct comparison between models.}\label{tab:auxiliary}
\gentable{tab_auxiliary}
\end{table}

\subsection{Exploratory stronger comparisons, information updates and reference scale}

\begin{table}[tbp]\centering\small
\caption{Exploratory stronger classical references with the latest input (German national analysis). The columns give WIS ratios [95\% intervals] for candidates with the latest input and with inputs from one and two weeks earlier. The original manuscript tabulated only the one-week-earlier column and showed the other two only in a figure; they are taken here from the same analysis output.}\label{tab:extended}
\gentable{tab_extended}
\end{table}

The comparison with a strengthened fixed ensemble and the selection by performance available in the past give different results (Table~\ref{tab:extended}). Neither is singled out to conclude ``superiority over classical methods in general'' or ``newer information is always better''. All comparisons, MAE, calibrated conditions and results by block length were saved.

Comparisons that separate updates of existing history from additions of new observation weeks in a fixed order, changes from a base week, and sensitivity to the reference version are shown in Appendix~\ref{app:updates}. Because the reference version changes the evaluation target itself, agreement with the latest version is not treated as biological correctness.

\subsection{Exploratory extensions to sites and another publication regime}\label{sec:explore}

In the exploratory extensions to German sites and New Zealand, no uniform relation was obtained. All comparisons by country and unit are shown in Appendix~\ref{app:panel}. New Zealand publishes roughly fortnightly, and many pairs with one-week-earlier input are identical to the latest input, so the magnitude of its ratios is not compared directly with Germany.

\subsection{Exploratory pooled learning, indicator definition and timing proxy}

Pooled learning gave different results for the German sites and for New Zealand. Details, including the sensitivity to the indicator definition and to the publication-time proxy, are given in Appendices~\ref{app:panel} and~\ref{app:sites}.

\subsection{Exploratory robustness to the age of input histories}

Relative robustness to older inputs is different from being able to replace forecasts made with the latest inputs. The check with a joint period and multiple comparisons could not confirm an advantage in robustness; details are in Appendices~\ref{app:age} and~\ref{app:joint}. The one-calendar-week difference in New Zealand is affected by fortnightly publication.

\subsection{Exploratory length of input histories}

Extending histories to at most 256 weeks did not yield a uniform improvement either. Some decreases of the ratio occurred in conditions in which the loss of the latest classical method increased; see Appendix~\ref{app:age}, which shows the losses of candidates and references separately.

\section{Discussion}

Surveillance staff who mainly use short horizons should not conclude, from the overall average that includes publication gaps, that the latest information can be replaced by general-purpose forecasts with one-week-earlier input. When a change of forecaster is considered, the difference between methods given the same information and the difference due to information updates within the same method should be placed side by side for the target pathogen. Combining more classical methods also changes the comparison, so a win or loss against a single weak reference should not be the basis for adoption. This study does not provide evidence for judging the cost-effectiveness of faster publication or model introduction, reductions in testing, or improvements in clinical alerts.

That a one-calendar-week-earlier input becomes identical to the latest input is itself a property that follows from the publication schedule. What should be regarded as findings is not its existence but the measured comparisons: the difference between the average that includes long horizons and the performance at short horizons, the differences between pathogens, and the change in how replaceability reads under stronger references. Differences across conditions result from simultaneous changes in models, epidemic periods and update processes, and are not interpreted as causal effects of a policy that increases observations.

The national comparisons depended not only on the type of model but also on the update state of the inputs, the choice of classical methods and the scale of the concentration referenced later. Comparisons that separate these factors can inform evaluation design when model updates and data updates are being considered. This study, however, did not measure the costs of either or the effects on field operations.

What to use as the latest classical forecast is also part of the comparison. A comparison only with the simple last value, or with a classical method selected after the fact as the best over the whole test period, does not adequately represent the actual options. This study kept both a fixed ensemble and a selection by past performance available at the time. Even so, it does not exhaust all forecasters or specialised wastewater models. Comparing the costs of model introduction and of faster publication would additionally require operational costs and decision utilities.

The choice of reference value also changes the property being measured. Matching the later level of concentration and capturing changes from the present are not the same. To avoid reading an ability to anticipate future revisions of the published normalised indicator as an ability to capture epidemic rises, both evaluations need to be shown side by side.

\section{Limitations}\label{sec:limitations}

First, this is a retrospective validation over publication histories, not forecasts actually submitted at the time. Git timestamps are a proxy for publication and do not independently guarantee the completeness of the history or the release times. Users at the time could have applied processing other than the eligibility rules chosen here. The origin and versions of the numbers were saved, but the immutability of an audit log has not been proven.

Second, the main reference is the published value after 56 days, not the true amount of pathogen or the true number of infected people. Changes in normalisation, participating sites, testing methods and sample counts can affect the values. Below-LOQ (limit of quantification) flags do not give true values either. Forward-filling of internal missing values turns part of the input into repeated values. Comparisons that exclude missing or out-of-range cases are limited to conditions with usable inputs and references.

Third, the number of epidemics in the validation period is limited. Multiple sites and overlapping two-week rises are not independent events, and varying the block length cannot handle all unknown dependence. The nominal intervals of the main and stratified tables are not adjusted for multiple comparisons. The added simultaneous intervals for the joint period cover only limited comparison families per country; they are neither an adjustment covering all results nor a strict coverage guarantee. Clinical alerts were not validated either.

Fourth, the pretrained models are limited to fixed public weights and a local implementation. The results are not generalised to new models in general or to future versions. No comparison was made in which the general-purpose models themselves are fine-tuned for the task. The pooled-learning comparison also changes features and training information, so the superiority of model architectures cannot be isolated. The past calibration forecasts of TimesFM, including those before its publication, are not a history of models that were available at the time; they must be interpreted separately from the uncorrected post-publication comparison.

Fifth, the sites, New Zealand, the additional classical methods and the information decomposition are exploratory checks planned after the main results. In topic selection, one critique of the topic candidates by a designated model was consulted. The model was \texttt{gpt-6-astra}, used on 2026-10-18, and the full critique and the response JSON were saved. Facts raised in the critique were checked separately, and those that could not be confirmed were not used as grounds. The exact model identifier of the research session could not be obtained,\footnote{The human author stated it after the run (\texttt{gpt-6.1-sol}; see the AI Declaration); it is not in the run's records.} so the difference from the critique model could not be checked directly from the runtime. The code, analyses, interpretation and paper were produced in this session and were not delegated to other models. The saved input instructions are an incomplete transcription, and their omissions are stated in the manifest and notes.

Sixth (added after the run), the run did not search the issue trackers of the data providers. A search of the AMELAG repository on 2026-10-19 found announcements by the maintainers concerning the test period: the addition of the flow-normalised variable from 20 August 2025~\cite{amelag36}, a change of smoothing method with renamed variables on 22 October 2025~\cite{amelag40}, and changed update days in calendar weeks 52 and 1 around the 2025/26 holidays~\cite{amelag45}. The run did not use these announcements, and whether they explain the update gaps described in Section~\ref{sec:data} was not examined. On the same day the head of the repository was still the frozen commit.

\section{Reproducibility and Research Process}

Plans made before the experiments, the reasons for changes and failed runs were all saved. The order of plan finalisation and the main runs can be traced from the local Git history and the event log, but Git times are self-reported and are not tamper-proof evidence. The correspondence between claims and evidence, and the recalculation from a single entry point with its run time, are documented in the replication package. Recalculation from saved forecasts requires no AI API, no weight downloads and no external network. The full replay for the initial draft, and the full re-run after the revision together with the final check of the paper and figures, are recorded in the replication package. Both are re-runs by the same researcher and not independent reproductions.

The review was a single round by another AI agent; the original text transmitted by a human and the decision on each point are saved in the replication package. The reviewer's model identifier and the scope of its re-runs are transmitted information that the research agent did not observe directly. The review is not called human peer review or independent reproduction.

% ---------------------------------------------------------------------------
\section*{AI Declaration}
\addcontentsline{toc}{section}{AI Declaration}

\paragraph{Role of AI in this work.} This study was conceived, designed, executed, analysed and written by an AI system acting autonomously: an OpenAI model running in the Codex command-line agent on a local Apple-silicon macOS workstation (macOS 26.5.2, 18\,GiB of memory, Python 3.13.7). The model was \texttt{gpt-6.1-sol} with reasoning effort xhigh, as stated by the human author after the run. The run's own records describe the model only as ``Codex, based on GPT-6'' (the developer description visible to the agent); the exact model identifier, the reasoning-effort setting and the runtime version were not exposed to the agent and are recorded in the run as unknown. The author's statement agrees with the process list of the machine, which after the run showed a Codex CLI process (installed package version 0.162.1) that had started 8\,s before the run with the options \texttt{-m gpt-6.1-sol -c model\_reasoning\_effort=xhigh} and the instruction protocol as its prompt. The AI selected the research topic itself from candidates it generated, searched and assessed prior work, wrote all software, collected and analysed all data, chose the statistical procedures, interpreted the results and drafted the original manuscript (in Japanese). The study was carried out in one agent session that started at 15:26:38 UTC on 18 October 2026; the initial draft was submitted at 03:54:02 UTC on 19 October, the review was received at 04:38:56 UTC, and the final manuscript was submitted at 06:38:53 UTC. The session stopped three times at the runtime's usage limit and was resumed by the human author each time. The pre-set budget was 12 hours of elapsed time until the initial draft and 2 hours for the review response (both excluding confirmed usage-limit stops), 10\,GiB of disk space, at most two concurrent experiment processes, 30 minutes per experiment process and USD 20 for AI API calls. Confirmed stops of 8,519\,s (research period) and 4,561\,s (review response) were excluded from the time budget; a second stop during the research period, whose length could not be measured, was not excluded. This English version was translated and formatted for arXiv by Claude Opus 5.5 (Anthropic; model identifier \texttt{claude-opus-5-5}) in Claude Code at the human author's request, after the research session had ended. In that work the same model re-checked every reference through the arXiv API, Crossref and the cited web pages, searched the issue tracker of the AMELAG data repository (which led to the sixth limitation in Section~\ref{sec:limitations}), regenerated all tables from the run's analysis outputs with the run's own rendering code, replaced the two figures of the original manuscript with tables, and re-ran the run's recalculation pipeline from the replication package, which reproduced all 70 analysis tables and the original manuscript exactly. \fillin{State whether the human author reviewed the translation.}

\paragraph{Role of the human author.} Ko Fujimura wrote a fixed, topic-agnostic instruction protocol for autonomous research (version 1.9), which defined resource limits, safety rules, permission procedures and record-keeping requirements, and started the session. The protocol left the choice of topic, method and interpretation to the AI, told it not to ask humans for advice on them or on improving the paper, and limited human involvement to permission and safety decisions, resuming execution and transmitting the result of one review round. During the run, the human author resumed the session three times after usage-limit stops and transmitted one review report; no permission requests were made, and no other input on scientific content was given. The review was written by a separate AI agent (Claude Opus 5.5 in Claude Code 2.1.303, according to that agent's own record), which worked under a fixed review protocol on a copy of the submitted initial draft. After the run, the human author requested this English version and the replication package. \fillin{State whether the human author has checked the data and results.}

\paragraph{Tools and data used by the AI.} The AI used the runtime's shell and file tools, a clock, and web search and page retrieval (the number of searches was not recorded). As the protocol allowed, it made one call to the OpenAI API (model \texttt{gpt-6-astra}, reasoning effort xhigh, 18 October 2026) to obtain a critique of five topic candidates; the recorded cost, calculated from the returned usage and the published price list and not reconciled with an invoice, was USD 2.2708675. No other AI model, sub-agent or paid API was used for the research. Data: the AMELAG wastewater surveillance data of the Robert Koch Institute and the German Environment Agency (CC BY 4.0); the New Zealand wastewater surveillance data published by PHF Science, formerly ESR (CC BY 4.0); and, for a pilot of a topic candidate that was not selected, annual air-quality data of the US Environmental Protection Agency. Models and software: the public weights of Chronos-Bolt base (Amazon, Apache-2.0) and TimesFM 2.5 200M (Google, Apache-2.0), run locally on CPU, with chronos-forecasting 2.3.2, timesfm 2.0.2, statsforecast 2.1.1, scikit-learn 1.9.1, PyTorch 2.14.1 and pandas 2.3.3. No human participants were involved, and no personal data were collected.

\paragraph{Verification and its limits.} Unless stated otherwise above, all verification was performed by the same AI that produced the work. This included an offline recalculation of all tables from the saved forecasts through a single entry point, run once for the initial draft (43 stages, 64 tables, 177,911 rows) and once after the revision (44 stages, 70 tables, 178,110 rows), each into an empty output directory and each with a maximum absolute numeric difference of 0.0 from the saved tables; checks of every claim against its evidence rows; a separate standard-library implementation of the temporal checks; and a separate implementation of the weekly site aggregation, compared over all 79 saved versions and 157 series (1,246,056 weekly values). The first attempt at the full recalculation failed because two quality tables processed datasets in a different order; it is kept in the record. The reviewing AI agent re-ran the recalculation in a newly created environment, regenerated all national and New Zealand forecasts with byte-identical results and recomputed the main ratios with its own code; this was another AI agent working on the same machine with the same weight files, not an independent replication. All of this is self-verification by AI systems and is not independent peer review or independent replication. Several analyses were decided after seeing results; the manuscript labels them as exploratory, and every plan revision is preserved with its timestamp.

\paragraph{Potential conflict of interest.} The AI that performed this study is developed by OpenAI, and the topic critique also used an OpenAI model. The study evaluates pretrained time-series models developed by Amazon (Chronos-Bolt) and Google (TimesFM) against classical statistical methods; no relation between OpenAI and the evaluated systems was identified. The review during the run and this English version were produced by Claude Opus 5.5, developed by Anthropic; no relation between Anthropic and the evaluated systems was identified. Because the same model reviewed the original manuscript and prepared this version, this version was not prepared independently of the reviewer. The human author declares no competing interests.

\paragraph{Provenance.} The complete research record is available at \fillin{repository URL}. It includes the instruction protocol as transcribed by the AI at the start of the run (an incomplete manual transcription whose known omissions are documented), the environment manifest, the time-stamped event log, the research log, every plan version and revision, all code, raw and derived outputs (or SHA-256 hashes where data could not be redistributed), and a SHA-256 manifest of the artefacts. Errors found during the run are documented in the log rather than removed.

\paragraph{Authorship and responsibility.} In accordance with arXiv's policy on generative AI tools and the COPE and ICMJE positions that AI tools cannot meet authorship criteria, the AI system is not listed as an author. The human author takes full responsibility for the content of this preprint, including any errors in it, irrespective of how the content was generated.

\section*{Data and Code Availability}
\addcontentsline{toc}{section}{Data and Code Availability}

All code, analysis outputs, the instruction protocol, plan versions, event log and SHA-256 manifest are available at \fillin{repository URL}. Code is released under the MIT License, and data and text under the Creative Commons Attribution 4.0 International licence (CC BY 4.0); third-party material keeps its original licence. The AMELAG and New Zealand data files are redistributed unmodified under their CC BY 4.0 licences, with attribution and licence notices; some were gzip-compressed by the run, which recorded their SHA-256 before and after compression. The analysis outputs derived from them are adaptations made in this study; their use does not imply endorsement by the data providers. Not redistributed are the model weights of Chronos-Bolt and TimesFM (Apache-2.0; about 1.8\,GB), which can be downloaded at the pinned revisions recorded with their SHA-256 hashes; the full texts and captured web pages of the literature the AI read, which are identified by URL and hash; and one intermediate output file of 139.7\,MB that exceeds common repository file-size limits and is recomputed by the package's reproduction pipeline. The package is a snapshot of the run directory at its final administrative commit; the run's local Git history is not included, but the initial draft is kept in the package. No numerical value was changed relative to the original manuscript: the tables of this version were regenerated from the run's analysis outputs with the run's own rendering code, which reproduced the original manuscript byte for byte. Because the analysis outputs are stored unrounded and rounded only once for display, no rounding corrections were needed. The figures of the original manuscript were replaced by tables (Table~\ref{tab:extended} gained two columns; the second figure repeated Tables~\ref{tab:primary}, \ref{tab:extended} and~\ref{tab:panel} and was omitted).

\begin{thebibliography}{99}\small
\raggedright

\bibitem{adiga2026} A.~Adiga, J.~Chou, G.~Kaur, A.~Warren, S.~Venkatramanan, B.~Espinoza, B.~Lewis, J.~Crow, et al. EpiFlow: A framework for improving the utility of wastewater signals for disease forecasting. arXiv:2608.06671v1, 2026.

\bibitem{ahmad2026} T.~Ahmad. Time-series foundation models that understand data revisions (VINTAGE-TS). arXiv:2609.28576v1, 2026.

\bibitem{alhassan2025} F.~Alhassan, H.~Karami, A.~Bleichrodt, J.~M. Hyman, I.~C.~H. Fung, R.~Luo, and G.~Chowell. COVID-19 forecasting from U.S. wastewater surveillance data: A retrospective multi-model study (2022--2024). arXiv:2512.01074v1, 2025.

\bibitem{amelagdata} Robert Koch-Institut and Umweltbundesamt. Abwassersurveillance AMELAG [data set]. Zenodo, 2026. \url{https://doi.org/10.5281/zenodo.23364951}. Citation recommended by the dataset's README at commit 6586326ea292f4bfa243d8444d1efc8f4fef772d.

\bibitem{amelagreadme} Robert Koch-Institut and Umweltbundesamt. Abwassersurveillance AMELAG: README (normalisation and nationwide aggregation), at commit 6586326ea292f4bfa243d8444d1efc8f4fef772d. \url{https://github.com/robert-koch-institut/Abwassersurveillance_AMELAG/blob/6586326ea292f4bfa243d8444d1efc8f4fef772d/Readme.md}.

\bibitem{amelag36} RKI Open Data Team. Erweiterung der Datei \texttt{amelag\_einzelstandorte.tsv} und \texttt{amelag\_aggregierte\_kurve.tsv} um Variable \texttt{viruslast\_normalisiert}. GitHub issue \#36, Abwassersurveillance AMELAG, 13 August 2025. \url{https://github.com/robert-koch-institut/Abwassersurveillance_AMELAG/issues/36}. (Found after the run.)

\bibitem{amelag40} RKI Open Data Team. \"Anderungen in den Dateien \texttt{amelag\_einzelstandorte.tsv} und \texttt{amelag\_aggregierte\_kurve.tsv} am 22.10.2025. GitHub issue \#40, Abwassersurveillance AMELAG, 17 October 2025. \url{https://github.com/robert-koch-institut/Abwassersurveillance_AMELAG/issues/40}. (Found after the run.)

\bibitem{amelag45} RKI Open Data Team. Vorgezogene Datenaktualisierungen in KW 52 und ausgesetzte Datenaktualisierung KW 1. GitHub issue \#45, Abwassersurveillance AMELAG, 17 December 2025. \url{https://github.com/robert-koch-institut/Abwassersurveillance_AMELAG/issues/45}. (Found after the run.)

\bibitem{bracher2021} J.~Bracher, E.~L. Ray, T.~Gneiting, and N.~G. Reich. Evaluating epidemic forecasts in an interval format. \emph{PLOS Computational Biology}, 17(2):e1008618, 2021. \url{https://doi.org/10.1371/journal.pcbi.1008618}

\bibitem{bracher2022} J.~Bracher, E.~L. Ray, T.~Gneiting, and N.~G. Reich. Correction: Evaluating epidemic forecasts in an interval format. \emph{PLOS Computational Biology}, 18(10):e1010592, 2022. \url{https://doi.org/10.1371/journal.pcbi.1010592}

\bibitem{cdcpipeline} CDCgov. wastewater-informed-covid-forecasting: official pipeline README and model definition. GitHub repository, \url{https://github.com/CDCgov/wastewater-informed-covid-forecasting} (no year stated; accessed 18 October 2026).

\bibitem{chronosbolt} Amazon. amazon/chronos-bolt-base, model card. Hugging Face, \url{https://huggingface.co/amazon/chronos-bolt-base}.

\bibitem{croushore2001} D.~Croushore and T.~Stark. A real-time data set for macroeconomists. \emph{Journal of Econometrics}, 105(1):111--130, 2001. \url{https://doi.org/10.1016/S0304-4076(01)00072-0}

\bibitem{giannone2005} D.~Giannone, L.~Reichlin, and D.~Small. Nowcasting GDP and inflation: The real-time informational content of macroeconomic data releases. Finance and Economics Discussion Series 2005-42, Federal Reserve Board, September 2005. \url{https://www.federalreserve.gov/pubs/feds/2005/200542/200542pap.pdf}

\bibitem{hewitt2022} J.~Hewitt, S.~Trowsdale, B.~A. Armstrong, J.~R. Chapman, K.~M. Carter, D.~M. Croucher, C.~R. Trent, R.~E. Sim, and B.~J. Gilpin. Sensitivity of wastewater-based epidemiology for detection of SARS-CoV-2 RNA in a low prevalence setting. \emph{Water Research}, 211:118032, 2022. \url{https://doi.org/10.1016/j.watres.2021.118032}

\bibitem{hu2026} Y.~Hu, J.~Han, and M.~Liu. Protocol-aware epidemic forecasting across heterogeneous public health surveillance systems. \emph{Frontiers in Public Health}, 14:1829302, 2026. \url{https://doi.org/10.3389/fpubh.2026.1829302}

\bibitem{fpp3} R.~J. Hyndman and G.~Athanasopoulos. \emph{Forecasting: Principles and Practice}, 3rd edition (chapters on exponential smoothing and ARIMA). OTexts, \url{https://otexts.com/fpp3/}.

\bibitem{kalahasti2025} S.~Kalahasti, B.~Faucher, B.~Wang, C.~Ascione, R.~Carbajal, M.~Enault, C.~V. Cassis, T.~Launay, et al. Foundation time series models for forecasting and policy evaluation in infectious disease epidemics. medRxiv, 2025. \url{https://doi.org/10.1101/2025.02.24.25322795} (Only the abstract was checked.)

\bibitem{lison2026} A.~Lison, R.~E. McLeod, J.~S. Huisman, J.~D. Munday, C.~Ort, T.~R. Julian, and T.~Stadler. Real-time estimation of pathogen transmission dynamics from wastewater. \emph{Nature Communications}, 17:8563, 2026. \url{https://doi.org/10.1038/s41467-026-75380-3}

\bibitem{mcdonald2021} D.~J. McDonald, J.~Bien, A.~Green, A.~J. Hu, N.~DeFries, S.~Hyun, N.~L. Oliveira, J.~Sharpnack, et al. Can auxiliary indicators improve COVID-19 forecasting and hotspot prediction? \emph{Proceedings of the National Academy of Sciences}, 118(51):e2111453118, 2021. \url{https://doi.org/10.1073/pnas.2111453118}

\bibitem{mcgough2020} S.~F. McGough, M.~A. Johansson, M.~Lipsitch, and N.~A. Menzies. Nowcasting by Bayesian smoothing: A flexible, generalizable model for real-time epidemic tracking. \emph{PLOS Computational Biology}, 16(4):e1007735, 2020. \url{https://doi.org/10.1371/journal.pcbi.1007735}

\bibitem{nzdocs} ESR-NZ. Aotearoa New Zealand Wastewater COVID-19 Surveillance Programme: official repository documentation. GitHub repository, \url{https://github.com/ESR-NZ/covid_in_wastewater}.

\bibitem{nzdata} COVID-19 Data Repository by the New Zealand Institute for Public Health and Forensic Science Ltd (PHF Science). \url{https://github.com/ESR-NZ/covid_in_wastewater}.

\bibitem{panja2026} M.~Panja, D.~D'Agostino, H.~Li, T.~Chakraborty, and N.~Liu. EpiCastBench: Datasets and benchmarks for multivariate epidemic forecasting. arXiv:2605.11598v1, 2026.

\bibitem{radermacher2025} J.~Radermacher, S.~Thiel, A.~Kannt, and H.~Fr\"ohlich. Wastewater as an early indicator for short-term forecasting COVID-19 hospitalization in Germany. \emph{BMC Public Health}, 25:2910, 2025. \url{https://doi.org/10.1186/s12889-025-24149-2}

\bibitem{romano2019} Y.~Romano, E.~Patterson, and E.~J. Cand\`es. Conformalized quantile regression. In \emph{Advances in Neural Information Processing Systems 32}, 2019. \url{https://papers.nips.cc/paper/2019/hash/5103c3584b063c431bd1268e9b5e76fb-Abstract.html}

\bibitem{sklearnhgb} scikit-learn developers. HistGradientBoostingRegressor, scikit-learn 1.9.1 documentation. \url{https://scikit-learn.org/stable/modules/generated/sklearn.ensemble.HistGradientBoostingRegressor.html}.

\bibitem{tang2025} J.~Tang, A.~Rumack, B.~Wilder, and R.~Rosenfeld. Real-time forecasting of data revisions in epidemic surveillance streams. \emph{PLOS Computational Biology}, 21(11):e1013709, 2025. \url{https://doi.org/10.1371/journal.pcbi.1013709}

\bibitem{timesfm25} Google. google/timesfm-2.5-200m-pytorch, model card. Hugging Face, \url{https://huggingface.co/google/timesfm-2.5-200m-pytorch}.

\end{thebibliography}

% ---------------------------------------------------------------------------
\clearpage
\appendix
\section{Supplementary Results}\label{app:results}

All calibration conditions, 4-, 8- and 12-week blocks, pathogens, target positions, 2026 only, reference versions and input relations are in the respective analysis outputs. Results for favourable conditions only are not substituted for the main results. Outputs for additional rise ratios, ROC, LOQ flags, missing and out-of-range inputs and the sensitivity to the timing proxy can also be reached from the reproduction procedure.

\subsection{Publication schedule and distribution of identical inputs (post-review)}\label{app:schedule}

\begin{table}[H]\centering\small
\caption{Gaps in the updates of the saved German national aggregate file.}\label{tab:pubgaps}
\gentable{tab_pub_gaps}
\end{table}

\begin{table}[H]\centering\small
\caption{Number of identical-input pairs per issue week (German national analysis). ``---'' means that there was no identical-input pair; before the publication of TimesFM, it is outside the evaluation. The evaluation periods differ between the models.}\label{tab:identityweeks}
\gentable{tab_identity_weeks}
\end{table}

\begin{table}[H]\centering\small
\caption{Distribution of the latest-reference horizon and of identical inputs (German national analysis).}\label{tab:horizondist}
\gentable{tab_horizon_distribution}
\end{table}

\begin{table}[H]\centering\small
\caption{Intervals between publication days in New Zealand within the test period. This aggregates the saved publication history, including versions later excluded by the eligibility rule; it is not the working frequency of laboratories.}\label{tab:nzintervals}
\gentable{tab_nz_intervals}
\end{table}

\subsection{All comparisons for sites, New Zealand and pooled learning}\label{app:panel}

\begin{table}[H]\centering\small
\caption{Exploratory comparisons with one-week-earlier input for German sites and New Zealand. The numbers of sites and regions are not numbers of independent epidemics. NZ national and regional series contain the same epidemics, so mixed values are not used as the main grounds for generalisation. These are not rankings of countries or sites: targets, scales, periods and missing inputs or references differ, so rows cannot be compared directly. New Zealand publishes roughly fortnightly, so a calendar one-week-earlier input includes both identical inputs and inputs that differ by several observation weeks; it does not mean the same observation difference as in Germany. All results, including strata for missing data and information updates, are saved in the corresponding analysis outputs.}\label{tab:panel}
\gentable{tab_panel}
\end{table}

\begin{table}[H]\centering\small
\caption{Comparison with pooled learning (exploratory). The pooled model was trained separately for each country. It additionally uses matured reference values from other series, the pathogen category and the annual cycle of the target date. It is not an architecture comparison that equalises training information with the univariate general-purpose models, but a check that adds an operationally available comparator.}\label{tab:pooled}
\gentable{tab_pooled}
\end{table}

Comparisons of the non-normalised SARS-CoV-2 indicator on common samples and the sensitivity to shifting the publication-time proxy are shown in Appendix~\ref{app:sites}. Including pooled learning, no particular reference is selected to replace the conclusions of the main analysis.

\subsection{Robustness to input age and history length}\label{app:age}

Table~\ref{tab:interaction} shows the group with an extended observation end in which all four forecasts (general-purpose and classical, each with old and latest input) exist for the same targets. The information age is one week; 56-day reference, uncorrected, full test period, after model publication.

\begin{table}[H]\centering\small
\caption{Relative loss increase from old to latest input for general-purpose and classical forecasts (exploratory). A ratio of ratios below 1 means a smaller relative loss increase for the general-purpose model. The four forecasts are compared together by resampling the same weeks; the interaction is not judged from the overlap of separate intervals. That the ratio of a fixed model equals 1 with identical input is a consistency check, not an empirical finding.}\label{tab:interaction}
\gentable{tab_interaction}
\end{table}

The ratio of ratios was small in the comparison with the German national fixed three-method ensemble, but changing to six methods or to New Zealand did not give uniform results (Table~\ref{tab:interaction}). All comparisons with past-performance selection and pooled learning were also saved. ``Old'' here refers to the cutoff of the numeric history; it is not a condition in which all information used for selection or training is rolled back. Past-performance selection can choose different components on the basis of past performance even when the current input is the same.

For TimesFM, even in the comparison with the national fixed three-method ensemble, the interval of the ratio of ratios includes 1 with 4-week blocks or for 2026 alone. Robustness over all periods and block lengths is not concluded.

In an exploratory sensitivity analysis that aligned the extended-observation-end group to the joint period after both models' publication and jointly resampled 32 statistics per country, the simultaneous intervals of the ratio of loss ratios (national, fixed three methods, one week earlier) were 0.882 [0.766, 1.017] for Chronos-Bolt and 0.940 [0.816, 1.083] for TimesFM. Both include 1, so an advantage in relative robustness is not confirmed from the nominal intervals alone. All comparison families and results by block length are shown in Appendix~\ref{app:joint}.

Inputs of at most 104 weeks and at most 256 weeks were compared exploratorily. The published versions, observation ends, target weeks, horizons and the 56-day reference are the same, and the comparison is limited to targets for which all four forecasts exist (short and long general-purpose, short and long latest classical). TimesFM was computed with the maximum context length of 256 of its inference setting for both short and long inputs, so ``both short'' is not necessarily identical in setting to the original main table. Tables~\ref{tab:history} and~\ref{tab:historylosses} use the latest fixed three-method ensemble as reference, with one-week-earlier input, uncorrected, full test period, after model publication and 8-week blocks.

\begin{table}[H]\centering\small
\caption{History length of at most 104 versus at most 256 weeks (exploratory). WIS ratios relative to the latest fixed three-method ensemble.}\label{tab:history}
\gentable{tab_history}
\end{table}

\begin{table}[H]\centering\small
\caption{Changes of the losses of the candidate and of the latest classical reference on the same targets when the history is extended (exploratory). Cases whose histories were not extended are included in the denominator, so the number of extended cases is shown in Table~\ref{tab:history}.}\label{tab:historylosses}
\gentable{tab_history_losses}
\end{table}

Agreement in conditions whose input was originally at most 104 weeks long and whose numeric columns are unchanged is a consistency check and is not counted as an effect of longer histories. All references, the extended group, the extended-observation-end group, 2026, MAE and 4-, 8- and 12-week blocks were saved.

In the comparison with the NZ national fixed three-method ensemble, the ratio when both are long was 1.052 [0.914, 1.219] for Chronos-Bolt and 1.035 [0.913, 1.165] for TimesFM. The long/short loss ratio of the latest classical method, however, was 1.115 [1.071, 1.167] and 1.111 [1.078, 1.183], respectively, so the loss on the reference side also increased. This is distinguished from the column in which only the general-purpose model is lengthened and the short latest reference is fixed. A decrease of the comparison ratio alone is not taken as evidence that the general-purpose model compensated better for old information.

Even when the same point estimate as in a main table is repeated, the interval endpoints differ slightly because each analysis used its own random stream for block resampling. The original saved values were not altered to make them agree. For the intervals of the main comparisons, refer to the main tables.

\subsection{Updates of existing history, observation additions and reference scale}\label{app:updates}

The intervals of identical statistics in this appendix also include endpoint differences caused by resampling independent of the main analysis. The agreement of point estimates and targets was checked, and interval differences are not read as reversals of results.

WIS measured as change from a base week was compared with WIS measured on levels (Table~\ref{tab:anchored}). Change from the base week is not a clinical ground truth either; it is a sensitivity analysis using a different measurement target.

\begin{table}[H]\centering\small
\caption{WIS on levels (``WIS'') and on changes from a common base week (``anchored WIS''), one-week-earlier input relative to the latest fixed three-method ensemble.}\label{tab:anchored}
\gentable{tab_anchored}
\end{table}

In the decomposition with a fixed update order, the loss difference A$-$B is described as the update of existing history and B$-$C as the addition of observations of new weeks (Table~\ref{tab:updatecontrib}). That the differences add up is an identity and is not a finding.

\begin{table}[H]\centering\small
\caption{Decomposition of the WIS reduction into the revision of existing history and the addition of new observation weeks.}\label{tab:updatecontrib}
\gentable{tab_update_contrib}
\end{table}

To avoid differences in the denominator caused by missing references, Table~\ref{tab:maturity} compares on the same targets for which all reference versions exist.

\begin{table}[H]\centering\small
\caption{WIS ratios by reference vintage on common targets (one-week-earlier input relative to the latest fixed three-method ensemble).}\label{tab:maturity}
\gentable{tab_maturity}
\end{table}

Table~\ref{tab:matrix} separates the difference between actually published input and input obtained by truncating the final version, from the difference in which reference is used. The final-version input contains future revisions. The differences in this matrix are not treated as the superiority or inferiority of available methods.

\begin{table}[H]\centering\small
\caption{Mean WIS by input history and reference (Chronos-Bolt with the latest input).}\label{tab:matrix}
\gentable{tab_input_reference_matrix}
\end{table}

\subsection{Measurement conditions of sites and publication regimes}\label{app:sites}

\begin{table}[H]\centering\small
\caption{German sites by stratum (exploratory). The share of below-LOQ flags is the mean of daily flags within the week; a share of 0 does not guarantee that all genes were quantifiable. According to the official description, values below the LOQ are replaced by half of the limit. Concentrations with flags are not treated as true viral loads or independent detection events. The strata also differ in epidemic period and other respects, so they are not causal comparisons between strata. Even with identical input the forecasters differ, so the forecasts themselves are not the same; the number of near-identical forecasts in the main comparison can be checked in the saved tables.}\label{tab:sitestrata}
\gentable{tab_site_strata}
\end{table}

For SARS-CoV-2 only, the non-normalised indicator and the normalised indicator of the main analysis were aligned to the same targets that can be evaluated with both (Table~\ref{tab:indicator}). The comparison is limited to the period in which an explicitly published non-normalised series is available.

\begin{table}[H]\centering\small
\caption{Non-normalised versus normalised SARS-CoV-2 indicator on common pairs (exploratory). Forecasts and reference values differ between the indicators. The last column is the difference of ratios with the same weeks resampled, not the causal effect of intervening on the normalisation.}\label{tab:indicator}
\gentable{tab_indicator}
\end{table}

\begin{table}[H]\centering\small
\caption{Sensitivity to the publication-time proxy (exploratory). A hypothetical condition in which the release of the inputs is uniformly delayed. The time of using past performance was also restricted to the later of the reference maturity cutoff and the delayed release time. This is not a result confirming that actual release delays had these lengths.}\label{tab:timing}
\gentable{tab_timing}
\end{table}

\subsection{Sensitivity to interval correction}\label{app:calibration}

\begin{table}[H]\centering\small
\caption{Moving-window interval correction (exploratory). All windows are kept, and the best window was not chosen on test performance. It was verified that the median and the MAE of point forecasts do not change on the common sample. The change in WIS from interval correction is not interpreted as an improvement of point forecasts.}\label{tab:calibration}
\gentable{tab_calibration}
\end{table}

\subsection{Missing inputs, constant series and imputation}\label{app:missing}

\begin{table}[H]\centering\small
\caption{Availability of one-week-earlier inputs. Even when an input is available, no loss pair can be formed if the later reference value or the comparator is missing. The pair counts of the main results are not to be confused with the total of available inputs.}\label{tab:availability}
\gentable{tab_availability}
\end{table}

\begin{table}[H]\centering\small
\caption{Audit of all saved unique numeric histories. This includes the calibration period and diagnostic inputs truncated from future national versions, so it is neither the frequency among test-period published inputs nor a count of independent cases. The provenance of measured versus imputed values is not determined by input IDs alone; the case-level audit below uses the records of each published version.}\label{tab:uniqueinputs}
\gentable{tab_unique_inputs}
\end{table}

\subsection{float32 conversion, close inputs and imputation conditions}\label{app:float32}

\begin{table}[H]\centering\small
\caption{Identity of one-week-earlier and latest inputs under several definitions. All test-period actually published inputs were compared without selecting on the availability of a later reference, so the denominators differ from the main loss tables. float32 identity is also limited to histories with the same dates and length. The difference thresholds do not imply biological equivalence.}\label{tab:float32}
\gentable{tab_float32}
\end{table}

\begin{table}[H]\centering\small
\caption{German sites by imputation condition (exploratory). The groups without imputation are small, and period, pathogen and observation frequency also vary between groups. Differences between strata are not interpreted as causal effects of imputation. All thresholds, all references, 2026, MAE and block lengths were saved.}\label{tab:imputation}
\gentable{tab_imputation}
\end{table}

\subsection{Additional references for long histories and numerical compatibility}\label{app:historyrefs}

Table~\ref{tab:historyrefs} (on a landscape page) repeats the history-length comparison with the other latest references.

\begin{sidewaystable}\centering\small
\caption{History length with further latest references (exploratory). The long/short loss ratios of the numerator and denominator for all references were also saved. The pooled model is a supplementary reference with the existing latest 104-week features and trained models kept fixed, so the history lengths of the two sides are not matched. Classical components that use a fixed recent window keep that window.}\label{tab:historyrefs}
\adjustbox{max width=\textheight}{%% generated by preprint/make_tables_tex.py from the run's analysis outputs; do not edit by hand
\begin{tabular}{lllrrrrr}
\toprule
Dataset & \hdl{Pretrained model\\(1 week earlier)} & \hdl{Latest-input\\reference} & \hdr{Common\\cases} & \hdr{Pretrained history\\extended} & \hdr{Both short:\\WIS ratio} & \hdr{Pretrained only long:\\WIS ratio} & \hdr{Both long:\\WIS ratio} \\
\midrule
Germany, national & Chronos-Bolt & Fixed 6-method & 338 & 186 & 1.127 [1.016, 1.223] & 1.150 [0.995, 1.292] & 1.134 [0.975, 1.282] \\
Germany, national & Chronos-Bolt & Past selection & 337 & 186 & 1.004 [0.927, 1.117] & 1.024 [0.906, 1.160] & 1.080 [0.952, 1.229] \\
Germany, national & TimesFM & Fixed 6-method & 262 & 160 & 1.102 [1.031, 1.190] & 1.132 [1.039, 1.231] & 1.118 [1.022, 1.225] \\
Germany, national & TimesFM & Past selection & 262 & 160 & 0.975 [0.872, 1.119] & 1.001 [0.894, 1.139] & 1.074 [0.983, 1.205] \\
Germany, sites & Chronos-Bolt & Fixed 6-method & 8043 & 4330 & 1.010 [0.953, 1.068] & 1.010 [0.949, 1.075] & 1.009 [0.944, 1.074] \\
Germany, sites & Chronos-Bolt & Past selection & 7751 & 4320 & 0.947 [0.873, 1.024] & 0.947 [0.874, 1.030] & 0.972 [0.899, 1.058] \\
Germany, sites & Chronos-Bolt & Pooled (history fixed) & 8043 & 4330 & 1.150 [1.044, 1.235] & 1.150 [1.055, 1.236] & 1.150 [1.055, 1.236] \\
Germany, sites & TimesFM & Fixed 6-method & 6216 & 3820 & 1.061 [0.993, 1.149] & 1.070 [1.005, 1.156] & 1.071 [1.000, 1.166] \\
Germany, sites & TimesFM & Past selection & 6206 & 3820 & 1.019 [0.937, 1.106] & 1.028 [0.949, 1.119] & 1.063 [0.986, 1.162] \\
Germany, sites & TimesFM & Pooled (history fixed) & 6216 & 3820 & 1.213 [1.116, 1.340] & 1.224 [1.125, 1.348] & 1.224 [1.125, 1.348] \\
NZ, national & Chronos-Bolt & Fixed 6-method & 114 & 114 & 1.240 [1.048, 1.533] & 1.154 [0.982, 1.383] & 0.891 [0.751, 1.072] \\
NZ, national & Chronos-Bolt & Past selection & 114 & 114 & 1.177 [0.974, 1.402] & 1.095 [0.946, 1.243] & 1.073 [0.932, 1.248] \\
NZ, national & Chronos-Bolt & Pooled (history fixed) & 114 & 114 & 0.907 [0.665, 1.182] & 0.844 [0.608, 1.075] & 0.844 [0.608, 1.075] \\
NZ, national & TimesFM & Fixed 6-method & 88 & 88 & 1.117 [0.955, 1.219] & 1.106 [0.999, 1.213] & 0.876 [0.774, 0.941] \\
NZ, national & TimesFM & Past selection & 88 & 88 & 1.081 [0.922, 1.246] & 1.070 [0.907, 1.300] & 1.065 [0.934, 1.203] \\
NZ, national & TimesFM & Pooled (history fixed) & 88 & 88 & 0.771 [0.607, 0.860] & 0.763 [0.637, 0.843] & 0.763 [0.637, 0.843] \\
NZ, regional & Chronos-Bolt & Fixed 6-method & 564 & 564 & 1.152 [1.069, 1.242] & 1.131 [1.059, 1.218] & 1.025 [0.945, 1.106] \\
NZ, regional & Chronos-Bolt & Past selection & 564 & 564 & 1.084 [0.982, 1.169] & 1.064 [0.994, 1.128] & 1.041 [0.968, 1.102] \\
NZ, regional & Chronos-Bolt & Pooled (history fixed) & 564 & 564 & 0.926 [0.811, 1.032] & 0.909 [0.796, 0.996] & 0.909 [0.796, 0.996] \\
NZ, regional & TimesFM & Fixed 6-method & 434 & 434 & 1.082 [1.022, 1.126] & 1.074 [1.012, 1.128] & 0.985 [0.919, 1.029] \\
NZ, regional & TimesFM & Past selection & 434 & 434 & 1.014 [0.932, 1.078] & 1.006 [0.953, 1.060] & 0.995 [0.941, 1.042] \\
NZ, regional & TimesFM & Pooled (history fixed) & 434 & 434 & 0.883 [0.773, 0.957] & 0.876 [0.791, 0.934] & 0.876 [0.791, 0.934] \\
\bottomrule
\end{tabular}
}
\end{sidewaystable}

The compatibility with the forecasts from the original short inputs was checked under the same conditions, including the calibration period (Table~\ref{tab:historycompat}).

\begin{table}[H]\centering\small
\caption{Numerical compatibility of the new short-input forecasts with the original forecasts. This is a numerical check by the same researcher, not an independent reproduction, and it includes differences in batching between models and in the TimesFM inference setting. The main results were not overwritten with the new forecasts. The maximum of 256 weeks is a fixed sensitivity condition and does not guarantee an optimal history length or the use of all available history.}\label{tab:historycompat}
\gentable{tab_history_compat}
\end{table}

\subsection{Exploratory sensitivity to the joint period and multiple comparisons}\label{app:joint}

This is an exploratory self-check that aligns both models to the same period after their publication. For each country, 16 comparisons were formed (2 models $\times$ 2 fixed references $\times$ 1 or 2 weeks earlier $\times$ 2 aggregation levels), and the ratio of ratios of each comparison together with the corresponding direct ratio gave 32 statistics. The same weeks were resampled jointly for all comparisons, and simultaneous intervals were formed from the 95th percentile of the maximum absolute deviation of the log ratios. There were 10000 resamples per country. This is neither an adjustment covering all comparisons in the paper nor a strict coverage guarantee for a small number of epidemics.

\begin{table}[H]\centering\small
\caption{Comparison families for the simultaneous intervals.}\label{tab:jointsummary}
\gentable{tab_joint_summary}
\end{table}

Table~\ref{tab:joint} shows the examples for the fixed three methods and one week earlier. The fixed six methods and two weeks earlier are also part of the simultaneous families, and all values were saved.

\begin{table}[H]\centering\small
\caption{Nominal and simultaneous intervals in the joint period (fixed three methods, one week earlier). Both intervals use the same period and the same forecasts. Differences from the original full-period Chronos-Bolt table include the effect of aligning the evaluation period. An interval that includes 1 is not taken as proof of no effect or of equivalence, and an exploratory interval that excludes 1 is not treated as an independent confirmation.}\label{tab:joint}
\gentable{tab_joint}
\end{table}

\section{Revision History of the Original Manuscript}

From the initial draft, the post-review analysis by horizon and the publication schedule were added, and the information age within the same method, the breakdown by pathogen and the breakdown of the few weeks were brought to the front. The main numbers and the claims of the initial draft were kept, and the interpretation of the overall average was narrowed. Exploratory tables were moved to the appendix, and figure captions and bibliographic entries were corrected. The grounds, changes and additional experiments for each review point are given in the review response in the replication package.

\end{document}
